\chapter{Jeux sous forme normale}
La théorie de la décision --par exemple du consommateur-- permet
de relier les choix d'un agent avec ses préférences. Un agent
rationnel est un agent qui cherche à maximiser son utilité
(espérée). Un ``bon'' choix, ou choix rationnel, est un choix
maximisateur d'utilité.

Nous allons introduire la théorie des jeux, que l'on pourrait
aussi appeler théorie de la décision interactive, en modélisant et
en étudiant des situations dans lesquelles plusieurs agents sont en
situations de choix, les choix des uns affectant l'utilité des
autres. Le ``bon'' choix pour un agent dépend maintenant des choix
des autres agents.

\section{Description d'un jeu sous forme normale}
\begin{exemple}
Amanda et Bertrand veulent se rencontrer à NY, mais n'ont pas
décidé d'un lieu de rendez-vous. Ils peuvent chacun se rendre soit
à l'Empire State Building, soit à Central Park.

Le but de chacun est de rencontrer l'autre, peu importe où. Les
préférences de chacun sur quel lieu o\`u aller dépendent donc de
ce que fait l'autre.

On peut donc décrire la situation par les éléments suivants :
\begin{itemize}
\item 2 joueurs, Amanda et Bertrand %
\item 2 décisions possibles pour chacun d'entre eux. (ESB, CP).%
\item L'utilité de chaque joueur en fonction de son choix et de celui
de l'autre joueur.
\end{itemize}
\end{exemple}

En général, un jeu sous forme normale est donné par
\begin{itemize}
\item  Une liste d'agents (ou joueurs) $I$; %
\item  Un ensemble de stratégies $S_i$ pour chaque agent $i$;%
\item  L'issue du jeu $s$ correspond au vecteur $(s_i)$ de
stratégies choisi; %
\item Des préférences pour les joueurs sur $S$,
représentées par une fonction $g_i \colon \Pi_i S_i\to \reals$.
\end{itemize}

\begin{exemple} Pour le jeu de rendez-vous à New-York,
\begin{itemize}
\item  $I=\{\mbox{Amanda},\mbox{ Bernard}\}$; %
\item  $S_i = \{E, C\}$; %
\item  $S = \{(E,C), (E,E), (C,C),
(C,E)\}$; %
\item $g_i(s_1, s_2) = 1$ si $s_1 = s_2$, $0$ sinon.
\end{itemize}
On représente le jeu par la matrice suivante : Amanda choisit la
ligne, et Bernard choisit la colonne. La paire de paiements pour
Amanda et pour Bernard est inscrite dans la case.


\begin{center}
\begin{game}{2}{2}
    &  $E$  &  $C$    \\
$E$ & $1,1$ & $0,0$   \\
$C$ & $0,0$ & $1,1$   \\
\end{game}\hspace*{20mm}%
\end{center}

\end{exemple}

\begin{definition}
Un jeu $G$ sous forme normale est donné par :
\begin{itemize}
\item Un ensemble fini de joueurs $I$ ; %
\item Des ensembles de
stratégies $S_i$ ; %
\item Des fonctions de paiement $g_i\colon
\Pi_i S_i\to \reals$.
\end{itemize}
On utilise les notations :\\
$S=\Pi_i S_i$ \\
$S_{-i} = \Pi_{j\neq i}S_j$.
\end{definition}

Nous passons maintenant en revue quelques jeux ou classes de jeux
particulièrement importants.

\subsection{Jeux à intérêts communs}
Ce sont des jeux dans lesquels tous les joueurs ont les mêmes
intérêts, ou préférences. On a donc $g_i = g_j$ pour tous joueurs
$i,j$. On peut voir ces jeux comme une généralisation des
problèmes à un agent.

Par exemple, le jeu de rendez-vous à New-York fait partie de cette
catégorie.

\subsection{Jeux à somme nulle}
Ce sont des jeux à deux joueurs dans lesquels les intérêts sont
parfaitement antagonistes. Ces jeux sont à l'opposé des jeux à
intérêts communs. Ce qui est gagné par un joueur est perdu par
l'autre. La somme des fonctions d'utilité est donc $0$. C'est à
dire que $g_i = -g_j$.

\begin{exemple} Jeux de pile ou face. Deux joueurs annoncent
Pile (P) ou face (F). Le joueurs 2 paie 1 \euro\ à 1 si les deux
choisissent la même chose, sinon le joueur 1 paie 1 \euro\ au
joueur 2.\\
\begin{center}
\begin{game}{2}{2}
    &  $P$  &  $F$    \\
$P$ & $1,-1$ & $-1,1$   \\
$F$ & $-1,1$ & $1,-1$   \\
\end{game}\hspace*{20mm}%
\end{center}
Chaque case représente l'utilité du joueur 1, suivie par celle du
joueur~2.
\end{exemple}

\subsection{La bataille des sexes}
Certains jeux font intervenir une part de coordination et une part de conflit entre
les agents. C'est le cas du jeu de la bataille des sexes suivant.\\


\begin{exemple}
Un couple veut décider d'une sortie.\\
L'homme préfère aller à l'Opéra, et
la femme préfère assister à un match de foot.\\
Pour chacun, être avec l'autre est plus important que le lieu.\\
\begin{center}
\begin{game}{2}{2}
    &  $F$  &  $O$    \\
$F$ & $2,1$ & $0,0$   \\
$O$ & $0,0$ & $1,2$   \\
\end{game}\hspace*{20mm}%
\end{center}
\end{exemple}

\subsection{La fureur de vivre} Deux adolescents en voiture foncent
l'un vers l'autre sur une route étroite. Personne ne veut sortir
de la route. Si les deux sortent, aucun n'est vraiment satisfait,
ni mécontent. Les choix pour chacun sont $F$ (faucon, rester sur
la route), ou $C$ (colombe, sortir de la route).
\begin{center}
\begin{game}{2}{2}
    &  $F$  &  $C$    \\
$F$ & $-1,-1$ & $10,0$   \\
$C$ & $0,10$ & $5,5$   \\
\end{game}\hspace*{20mm}%
\end{center}

\subsection{Dilemme du prisonnier} Deux prisonniers complices sont
interrogés séparément. Chacun peut trahir son partenaire en le
dénonçant ($T$), ou bien rester silencieux ($S$). Si les deux
trahissent, ils vont en prison pour 5 ans chacun. Si l'un trahit
et pas l'autre celui qui trahit sort libre et l'autre va en prison
pour 10 ans. Si personne ne trahit, ils vont en prison pour 3 ans
tous deux.
\begin{center}
\begin{game}{2}{2}
    &  $S$  &  $T$    \\
$S$ & $-3,-3$ & $-10,0$   \\
$T$ & $0,-10$ & $-5,-5$   \\
\end{game}\hspace*{20mm}%
\end{center}

Le jeu du dilemme du prisonnier est un exemple fondamental en
économie. Ce jeu est important car il fait ressortir une tension
entre l'intérêt individuel et l'intérêt collectif. De nombreuses
situations présentent une structure similaire à celle du dilemme
du prisonnier
\begin{itemize}
\item Achat par internet ; Un acheteur et un vendeur se sont mis
d'accord sur internet. Chacun a le choix entre envoyer le colis
(ou bien l'argent), et ne pas l'envoyer. Chacun préfère que
l'autre l'envoie et préférerait ne pas envoyer sa part. Cependant,
les deux sont plus satisfaits si chacun respecte sa
part du contrat plutôt que si aucun ne le fait.%
\item Course aux armements ; Chaque pays peut décider de s'armer,
ou non. L'intérêt des deux pays est qu'aucun ne dépense de
resources pour s'armer, mais à stratégie fixée de l'autre, chacun
préfère s'armer.%
\item Achat d'un 4x4 ; Un 4x4 est avantageux pour celui qui l'a car on peut impressionner
les autres voitures, et on se sent plus en sécurité. Mais ceci est au détriment des autres voitures. %
\item Collusion et la commission européenne ; La commission
européenne cherche à lutter contre les ententes secrètes des
industries. Le problème étant que personne ne veut dénoncer ces
ententes. La règle est que celui qui dénonce une entente ne se
verra pas poursuivi. Chacun a donc intérêt à dénoncer plutôt qu'à
être dénoncé.
\end{itemize}


\subsection{Compétition en quantités dite de Cournot}
Deux firmes, 1 et 2, produisent des biens identiques.\\
Chacune décide d'un niveau de production $q_i$, à un co\^ut
$c_i(q_i)$. Le prix résultant de la loi de l'offre et de la
demande est $p(q_1+q_2)$.
\begin{eqnarray*}
I   &=& \{1,2\}\\
S_i &=& \reals_+ \\
g_1(q_1,q_2) &=& q_1 p(q_1+q_2)-c_1(q_1)\\
g_2(q_1,q_2) &=& q_2 p(q_1+q_2)-c_2(q_2)
\end{eqnarray*}

\subsection{Compétition en prix dite de Bertrand}
Il s'agit d'un modèle de compétition par les prix. Chaque firme
décide d'un prix de vente du bien, et les consommateurs (une unité
en tout) achètent à la firme au prix le plus bas.
\begin{eqnarray*}
I   &=& \{1,2\}\\
S_i &=& [0,M] \\
\end{eqnarray*}
$$
g_i(p_i,p_j)=\begin{cases}
  p_i &\mbox{si }\;\; p_i < p_j\\
  0 &\mbox{si }\;\; p_i > p_j\\
  p_i/2 &\mbox{si }\;\; p_i = p_j
\end{cases}$$\\

\section{Stratégies dominantes et dominées}
Nous commençons maintenant l'étude des prédictions des issues des
jeux sous forme normale. Étant donné un jeu, et en supposant les
joueurs rationnels, que peut-on prédire qu'ils vont jouer, ou que
peut-on prédire qu'ils ne vont pas jouer ? Si nous devions jouer
dans un tel jeu, quel choix ferions-nous ?

Une stratégie dominée est une stratégie qui donne un paiement
strictement moins bon que celui d'une autre stratégie donnée, ceci
pour tout choix possible des adversaires.

\begin{definition}
Une stratégie $s_i$ est (strictement) dominée s'il existe $s'_i$
telle que
$$
\forall s_{-i}\in S_{-i},\, g_i(s'_i,s_{-i}) > g_i(s_i,s_{-i})
$$
On dit alors que $s'_i$ domine (strictement) $s_i$.
\end{definition}

Un joueur ``rationnel'' ne devrait jamais jouer une stratégie
dominée. En effet, il a la possibilité d'une autre stratégie dont
il sait qu'elle peut lui donner un meilleur gain, indépendamment
des choix des autres joueurs.

Une stratégie dominante donne un meilleur paiement que toute autre
stratégie, pour tous les choix des autres joueurs.

\begin{definition}
Une stratégie $s_i$ est (strictement) dominante si pour toute
autre stratégie $s'_i$.
$$
\forall s_{-i}\in S_{-i},\, g_i(s_i,s_{-i}) > g_i(s'_i,s_{-i})
$$
\end{definition}

La définition se réécrit :
\begin{proposition} Une stratégie $s_i$ est dominante si et seulement si elle domine toute stratégie
$s'_i\neq s_i$.
\end{proposition}

Une conséquence quasi-immédiate de la définition est : %
\begin{proposition}
Si une stratégie $s_i$ est dominante, alors elle est unique à
avoir cette propriété.
\end{proposition}

Si un joueur a une stratégie dominante, on peut alors penser que
cette stratégie constitue un bon choix. Pour chaque choix des
autres, elle donne le meilleur paiement possible.

Nous pouvons maintenant construire un raisonnement en se basant
sur le fait que si aucun joueur ne choisit de stratégie dominée,
alors les autres joueurs peuvent anticiper ce phénomène.

\begin{exemple}
Considérons le jeu suivant :
\begin{center}
\begin{game}{2}{2}
    &  $S$  &  $T$    \\
$S$ & $0,-2$ & $-10,-1$   \\
$T$ & $-1,-10$ & $-5,-5$   \\
\end{game}\hspace*{20mm}%
\end{center}

Que devrait jouer le joueur 1 s'il sait que le joueur 2 est
rationnel ?
\end{exemple}

\begin{exemple}
Considérons le jeu suivant :
\begin{center}
\begin{game}{2}{3}
    &  $G$  &  $M$   & $D$    \\
$H$ & $2,2$ & $1,1$  & $4,0$  \\
$B$ & $1,2$ & $4,1$  & $3,5$  \\
\end{game}\hspace*{20mm}%
\end{center}

\begin{itemize}
\item Que devrait jouer le joueur 1 s'il sait que le joueur 2 est rationnel ?\\
\item Que devrait jouer le joueur 2 s'il sait que le joueur 1 sait qu'il est rationnel ?\\
\end{itemize}

\end{exemple}

Ceci nous amène à définir le processus d'élimination itérée des stratégies dominées.

\begin{definition} La procédure d'élimination itérée des
stratégies dominées est la suivante :
\begin{itemize}
\item {\bf \'Etape 1 :} $S_i^0 = S_i$.%
\item {\bf \'Etape 2 :} $S_i^1 = \{$ stratégies non dominées
de $i$ $\}$.%
\item {\bf \'Etape $k+1$ :} $$S_i^{k+1} = \{s_i \in S_i^k,
\not\exists s'_i \in S^{k}_i, \forall s_{-i}\in
S^{k}_{-i}, g_i(s'_i,s_{-i})>g_i(s_i,s_{-i})\}$$%
Stratégies non dominées {\em face aux stratégies de } $S^k_i$%
\item {\bf \'Etape $\infty$ :} $S_i^{\infty} = \cap_k S^k_i$.
\end{itemize}
\end{definition}

Pour certains jeux, cette définition nous conduit à une prédiction
unique des stratégies suivies par les joueurs.

\begin{definition}
Un jeu $G$ est résoluble par élimination itérée de stratégies
dominées si pour tout $i$, $S_i^{\infty}$ est un singleton.
\end{definition}

\begin{exemple}
Que dit la procédure d'élimination de stratégies dominées dans le
jeu suivant ?
\begin{center}
\begin{game}{4}{4}
    &  $A$  &  $B$   & $C$   & $D$  \\
$A$ & $5,2$ & $2,6$  & $1,4$ & $0,4$\\
$B$ & $0,0$ & $3,2$  & $2,1$ & $1,1$\\
$C$ & $7,0$ & $2,2$  & $1,5$ & $5,1$\\
$D$ & $9,5$ & $1,3$  & $0,2$ & $4,8$\\
\end{game}\hspace*{20mm}%
\end{center}
\end{exemple}


Passons maintenant à l'exemple suivant.
\begin{exemple}~\\
\begin{center}
\begin{game}{2}{2}
    &  $G$  &  $D$\\
$H$ & $2,2$ & $4,4$\\
$B$ & $3,1$ & $4,1$\\
\end{game}\hspace*{20mm}%
\end{center}
Quelles sont les stratégies dominées, dominantes ?\\
Quelles stratégies paraissent ``raisonnables'' ?%
\end{exemple}

Dans l'exemple précédent, il n'y a pas de stratégies (strictement)
dominantes, ni dominées. Cependant, le choix de la stratégie $B$
pour le joueur 1, et $D$ pour le joueur 2, semble un bon choix,
car ces stratégies ne peuvent donner qu'un meilleur paiement que
l'autre, et jamais un paiement moins bon. Ceci nous conduit à la
définition d'une stratégie faiblement dominée.

\begin{definition} Une stratégie $s_i$ est faiblement
dominée s'il existe $s'_i$ telle que:
\begin{eqnarray*}
\forall s_{-i}, g_i(s'_i,s_{-i}) &\geq& g_i(s_i,s_{-i})\\
\exists s_{-i}, g_i(s'_i,s_{-i}) &>& g_i(s_i,s_{-i})
\end{eqnarray*}
On dit dans ce cas que $s'_i$ domine faiblement $s_i$.
\end{definition}

Dans le jeu précédent, $H$ est faiblement dominée pour le joueur
1, et $G$ est faiblement dominée pour le joueur 2.

\begin{definition}
Une stratégie $s_i$ est faiblement dominante si pour toute autre
stratégie $s'_i$.
\begin{eqnarray*}
\forall s_{-i}\in S_{-i},\, g_i(s_i,s_{-i}) &\geq &g_i(s'_i,s_{-i})\\
\exists s_{-i}\in S_{-i},\, g_i(s_i,s_{-i}) &> &g_i(s'_i,s_{-i})
\end{eqnarray*}
\end{definition}

Si une telle stratégie existe, elle est unique. Elle semble
représenter un choix ``raisonnable'', car on ne peut jamais le
regretter.

La définition se réécrit :
\begin{proposition} Une stratégie $s_i$ est faiblement dominante si et seulement
si elle domine faiblement toute stratégie $s'_i\neq s_i$.
\end{proposition}

Une conséquence quasi-immédiate de la définition est : %
\begin{proposition}
Si une stratégie $s_i$ est faiblement dominante, alors elle est
unique à avoir cette propriété.
\end{proposition}

\begin{exemple}
Un groupe de joueurs $I$ participe à une enchère pour un objet. La
valeur de cet objet est $v_i$ pour le joueur $i$. Si chaque joueur
$i$ met une enchère de $e_i$, l'objet est attribué à celui qui met
l'enchère la plus haute, à un prix correspondant à la seconde
enchère la plus haute.

On demande de montrer comme exercice que pour le joueur $i$, jouer
$e_i = v_i$ est une stratégie faiblement dominante.
\end{exemple}

Lorsqu'un joueur possède une stratégie dominante, on peut penser
que cette stratégie constitue un ``bon'' choix. En revanche, une
stratégie dominée (même faiblement) ne semble pas constituer un
bon choix. Nous sommes tentés de généraliser la procédure de
suppression itérée de stratégies dominées aux stratégies
faiblement dominées. Voyons ce que l'on obtient dans l'exemple
suivant :

\begin{exemple} Soit le jeux donné par la matrice de paiements

\begin{center}
\begin{game}{3}{2}
    &  $G$  &  $D$\\
$H$ & $1,1$ & $0,0$\\
$M$ & $1,1$ & $2,1$\\
$B$ & $0,0$ & $2,1$\\
\end{game}\hspace*{20mm}%
\end{center}

La procédure d'élimination des stratégies faiblement dominées nous
permet d'éliminer $T$, puis $L$. Cette procédure pourrait aussi
bien nous conduire à éliminer $B$, puis $R$.  Ces deux ordres de
suppression de stratégies dominées nous conduit à des résultats
différents. La méthode n'est donc pas concluante.
\end{exemple}

On demande de montrer en exercice que l'ordre de suppression des
stratégies strictement dominées n'a pas d'influence sur le
résultat final.

En conclusion : Si on supprime de manière itérée les stratégies
{\bf faiblement} dominées, le résultat peut dépendre de l'ordre
des itérations. Ce n'est donc pas une bonne procédure. En
revanche, nous n'avons pas ces difficultés avec la suppression
itérée de stratégies strictement dominées.

\begin{exemple}
Considérons le jeu de compétition de Cournot suivant :
\begin{itemize}
\item $c_i(q_i) = c q_i$. \item $p = \alpha - \beta(q_1+q_2)$.
\end{itemize}
Si $i$ croit que $j$ va choisir $q_j$, alors la meilleure réponse
(la stratégie qui maximise le paiement de $i$ à stratégie de $j$
fixée) est :
$$
MR_i(q_j)=
\begin{cases}
  \frac{\alpha-c}{2\beta} - \frac{q_j}2 &\mbox{si }\;\; q_i \leq\frac{\alpha-c}{\beta}\\
   0 &\mbox{sinon }
\end{cases}$$\\
Les joueurs choisissent dans $S^0=\reals$.\\
Toutes les stratégies hors de $S^1=[0,\frac{\alpha-c}{2\beta}]$
sont dominées
(par 0).\\
Ensuite, toutes les stratégies hors de
$S^2=[MR(\frac{\alpha-c}{2\beta}), MR(0)]$ sont dominées (par
$MR(\frac{\alpha-c}{2\beta})$). En itérant, on démontre que le jeu
est résoluble par élimination itérée de stratégies dominées et
que, $S_i^{\infty}=\{\frac{\alpha-c}{3\beta}\}$.
\end{exemple}

\section{\'Equilibre de Nash}
La suppression itérée de stratégies dominées est attractive car
elle repose sur l'idée de rationalité (et sur la connaissance par
les joueurs de la rationalité des autres, etc...). Cependant, peu
de jeux sont résolubles par élimination itérée de stratégies
dominées. Nous introduisons un concept qui s'applique à une plus
grande gamme de jeux : l'équilibre de Nash.

L'équilibre de Nash peut être compris comme une convention sociale
stable. C'est une règle collective de laquelle aucun individu n'a
intérêt à s'écarter pourvu que les autres individus respectent eux
aussi la règle.

\begin{definition} Un profil de stratégies $s^*=(s_i)_i$ est
un équilibre de Nash (en stratégies pures) si :
$$
\forall i, \forall s'_i\in S_i, g_i(s_i,s_{-i})\geq
g_i(s'_i,s_{-i})
$$
\end{definition}

La notion suivante d'équilibre de Nash strict est plus forte. Elle
suppose que les joueurs préfèrent tous strictement jouer
l'équilibre de Nash plutôt qu'une autre stratégie si les autres
joueurs suivent l'équilibre de Nash.
\begin{definition}  Un profil de stratégies $s^*\in S$ est un
équilibre de Nash {\em strict} (en stratégies pures) si
$$
\forall i, \forall s'_i \in S_i, g_i(s^*_i,s^*_{-i})>
g_i(s'_i,s^*_{-i})
$$
\end{definition}


En particulier, un équilibre de Nash strict est un équilibre de
Nash.

\begin{exemple} Dilemme du prisonnier
\begin{center}
\begin{game}{2}{2}
    &  $S$  &  $T$\\
$S$ & $-3,-3$ & $-10,-10$\\
$T$ & $0,-10$ & $-5,-5$\\
\end{game}\hspace*{20mm}%
\end{center}
Seul équilibre de Nash: $(T,T)$.
\end{exemple}

\begin{exemple} Soit le jeu
\begin{center}
\begin{game}{2}{3}
    &  $G$  &  $M$  &  $D$\\
$H$ & $2,2$ & $1,1$ & $4,0$\\
$B$ & $1,2$ & $4,1$ & $3,5$\\
\end{game}\hspace*{20mm}%
\end{center}
Unique équilibre de Nash : $(H,G)$.
\end{exemple}

Le jeu du dilemme du prisonnier est résoluble en stratégies
dominées, et l'équilibre de Nash est compatible avec cette
prédiction.

Le jeu précédent n'est pas résoluble en stratégies dominées, et
cependant l'équilibre de Nash nous donne une prédiction.

A l'équilibre de Nash, aucune stratégie dominée n'est jouée. Plus
généralement, toutes les stratégies des équilibres de Nash
survivent à l'élimination itérée des stratégies dominées.

\begin{proposition}
Si $s^*$ est un équilibre de Nash, alors pour tout $i$, $s^*_i\in
S^\infty_i$.
\end{proposition}
{\em Démonstration :} Supposons $s\in \Pi_i S^k_i$, $s_i\not\in
S^{k+1}_i$. Alors il existe $s'_i$ meilleur pour $i$ que $s_i$
contre tous $s_{-i}\in \Pi_{j\neq i} S^k_j$, donc en particulier
contre $s^*_{-i}$. Contradiction.

Si on a un profil de stratégies dominantes, ce profil constitue un
équilibre de Nash:
\begin{proposition}
Si $s^*$ est un profil de stratégies faiblement dominantes, $s^*$
est un équilibre de Nash. Si $s^*$ est un profil de strictement
dominantes alors $s^*$ est l'unique
équilibre de Nash. %
\end{proposition}

\begin{exemple} Le jeu suivant est-il résoluble
par élimination itérée de stratégies dominées ?\\
Quels sont ses équilibres de Nash ?

\begin{center}
\begin{game}{3}{3}
    &  $G$  &  $M$  &  $D$ \\
$h$ & $5,3$ & $0,4$ & $3,5$\\
$m$ & $4,0$ & $5,5$ & $4,0$\\
$b$ & $3,5$ & $0,4$ & $5,3$\\
\end{game}\hspace*{20mm}%
\end{center}
\end{exemple}

\begin{exemple}
Reprenons le jeu de compétition de Cournot donné par :
\begin{itemize}
\item $c_i(q_i) = c q_i$. \item $p = \alpha - \beta(q_1+q_2)$.
\end{itemize}
Quels sont les équilibres de Nash ?
\end{exemple}

\begin{exemple}
Deux firmes, 1, 2. Firme $i$ décide d'un prix $p_i$. Le co\^ut
marginal de production est $c$. La demande est donnée par $D(p)$,
tous les acheteurs achêtent au meilleur prix.
$$
D_i(p_1,p_2)=
\begin{cases}
  D(p_i)   &\mbox{si }\;\; p_i < p_j\\
  D(p_i)/2 &\mbox{si }\;\; p_i = p_j\\
  0        &\mbox{si }\;\; p_i > p_j
\end{cases}$$\\
On suppose $D$ décroissante, et $D(c)>0$.\\

Quels sont les équilibres de Nash ?
\end{exemple}

\begin{exemple}
Modèle de compétition spatiale, appelée aussi compétition de Hotelling.\\
Deux vendeurs de glace sur une plage linéaire ([0,1]).\\
Chacun choisit un emplacement $x_i\in [0,1]$. Les consommateurs
sont uniforméments répartis sur la plage,
et achètent au vendeur le plus proche.\\
Chaque vendeur veut maximiser ses ventes.\\
Représentation en jeux ?\\
\'Equilibres de Nash ?

\end{exemple}

Un jeu peut avoir plusieurs équilibres de
Nash.

\begin{exemple} Quels sont les équilibres de Nash en stratégies pures
du jeu du RdV à NY ?\\
\begin{center}
\begin{game}{2}{2}
    &  $E$  &  $C$  \\
$E$ & $1,1$ & $0,0$ \\
$C$ & $0,0$ & $1,1$ \\
\end{game}\hspace*{20mm}%
\end{center}
\end{exemple}

On peut aussi ne pas avoir d'équilibres. Voir le jeu suivant.
\begin{exemple}
Quels sont les équilibres de Nash en stratégies pures
du jeu suivant à somme nulle ?\\
\begin{center}
\begin{game}{2}{2}
    &  $P$  &  $F$  \\
$P$ & $-1,1$ & $1,-1$\\
$F$ & $1,-1$ & $-1,1$\\
\end{game}\hspace*{20mm}%
\end{center}
\end{exemple}

On peut avoir plusieurs équilibres de Nash, et cependant en
préférer un plutôt qu'un autre pour des raisons liées au
``risque'' pris en jouant les stratégies.

\begin{exemple}
Quels sont les équilibres de Nash en stratégies pures du jeu
suivant ?\\ Quelle stratégie choisiriez-vous ?
\begin{center}
\begin{game}{2}{2}
    &  $A$  &  $B$  \\
$A$ & $9,9$ & $-15,8$\\
$B$ & $8,-15$ & $7,7$\\
\end{game}\hspace*{20mm}%
\end{center}
\end{exemple}

Certaines stratégies faiblement dominées peuvent être jouées à
l'équilibre de Nash.

\begin{exemple}
Quels sont les équilibres de Nash en stratégies pures du jeu
suivant ?\\ Quelle stratégie choisiriez-vous ?
\begin{center}
\begin{game}{2}{2}
    &  $A$  &  $B$  \\
$A$ & $2,2$ & $1,2$\\
$B$ & $2,1$ & $1,1$\\
\end{game}\hspace*{20mm}%
\end{center}
\end{exemple}

\subsection{Interprétations de l'\'Equilibre de Nash}
La notion d'équilibre de Nash est centrale en théorie des jeux et
en microéconomie. Elle est aussi appliquée en informatique,
biologie, en sociologie... Les interprétations de cette notion
sont diverses. On peut citer les justifications suivantes :
\begin{itemize} %
\item Prescriptions. Un individu extérieur
propose un profil de stratégies. Si ce profil est un équilibre
de Nash, personne n'a intérêt à dévier.%
\item Communication. Les joueurs se mettent d'accord à l'avance
sur quel profil jouer. Les EN sont des accords viables.%
\item Introspection. Un joueur analyse le jeu et se convainc que
seul l'EN est possible. Valable uniquement si unique EN. %
\item Norme sociale. Dans certains pays, nous conduisons à
droite de la route, dans d'autres à gauche. C'est la norme sociale
qui dicte une règle. EN est nécessaire pour que cette
règle soit suivie.%
\item Apprentissage. Les joueurs découvrent les stratégies
suivies par les autres au cours d'un processus d'apprentissage.%
\item Evolution.  Si chaque espèce évolue de fa\c{c}on
adaptée aux autres espèces, la résultante est un EN.
\end{itemize}
